\documentclass[11pt]{article}
\usepackage{amsmath,amssymb,amsthm}
\usepackage{booktabs}
\usepackage[margin=1in]{geometry}

\newtheorem{theorem}{Theorem}
\newtheorem{lemma}[theorem]{Lemma}
\newtheorem{corollary}[theorem]{Corollary}
\newtheorem{proposition}[theorem]{Proposition}

\title{Problem 951: Euler Product Approximation}
\author{Project Euler Solutions}
\date{}

\begin{document}
\maketitle

\section{Problem Statement}

The Euler product for the Riemann zeta function is
\[
\zeta(s) = \prod_{p \text{ prime}} \frac{1}{1 - p^{-s}}.
\]
Let $E(s, N) = \prod_{p \le N} (1 - p^{-s})^{-1}$. Find $\lfloor E(3, 10^6) \times 10^{15} \rfloor$.

\section{Mathematical Foundation}

\begin{theorem}[Euler Product Formula]
For $\operatorname{Re}(s) > 1$,
\[
\zeta(s) = \sum_{n=1}^{\infty} n^{-s} = \prod_{p \text{ prime}} \frac{1}{1 - p^{-s}}.
\]
\end{theorem}

\begin{proof}
Each factor expands as a geometric series: $(1-p^{-s})^{-1} = \sum_{k=0}^{\infty} p^{-ks}$, convergent for $\operatorname{Re}(s) > 0$. For any finite set of primes~$P$,
\[
\prod_{p \in P} \sum_{k=0}^{\infty} p^{-ks} = \sum_{\substack{n \ge 1 \\ p \mid n \Rightarrow p \in P}} n^{-s}
\]
by the Fundamental Theorem of Arithmetic: every positive integer factors uniquely into primes, so expanding the product yields exactly one term~$n^{-s}$ for each~$n$ whose prime factors all lie in~$P$. As $P$ exhausts all primes, the right-hand side converges to~$\zeta(s)$ by absolute convergence of the Dirichlet series for $\operatorname{Re}(s) > 1$.
\end{proof}

\begin{lemma}[Monotonicity]
$E(s, N) \le \zeta(s)$ for all~$N$, with $E(s, N) \nearrow \zeta(s)$ as $N \to \infty$.
\end{lemma}

\begin{proof}
Each factor $(1-p^{-s})^{-1} > 1$ for $s > 0$. The partial product~$E(s,N)$ accounts for all positive integers whose prime factors are at most~$N$, a subset of all positive integers. Including additional primes strictly increases the product.
\end{proof}

\begin{theorem}[Error Bound]
\label{thm:error}
For $s = 3$ and $N \ge 2$,
\[
\zeta(3) - E(3, N) = O\!\left(\frac{1}{N^2 \ln N}\right).
\]
\end{theorem}

\begin{proof}
We have $\zeta(3)/E(3,N) = \prod_{p>N}(1-p^{-3})^{-1}$. Taking logarithms:
\[
\ln\frac{\zeta(3)}{E(3,N)} = \sum_{p>N} \sum_{k=1}^{\infty} \frac{1}{kp^{3k}} = \sum_{p>N} p^{-3} + O\!\Bigl(\sum_{p>N} p^{-6}\Bigr).
\]
By partial summation with the Prime Number Theorem $\pi(x) \sim x/\ln x$:
\[
\sum_{p>N} p^{-3} = -\frac{\pi(N)}{N^3} + 3\int_N^{\infty} \frac{\pi(t)}{t^4}\,dt \sim 3\int_N^{\infty} \frac{dt}{t^3 \ln t} \sim \frac{3}{2N^2 \ln N}.
\]
Since $\ln(1+\varepsilon) \sim \varepsilon$ for small~$\varepsilon$, we get $\zeta(3)/E(3,N)-1 = O(1/(N^2\ln N))$, and therefore $\zeta(3) - E(3,N) = O(1/(N^2\ln N))$.
\end{proof}

\begin{lemma}[Numerical Precision]
For $N = 10^6$, the relative error is approximately $3 \times 10^{-14}$. Standard IEEE 754 double-precision arithmetic suffices for extracting $\lfloor E(3,10^6) \times 10^{15} \rfloor$.
\end{lemma}

\begin{proof}
By Theorem~\ref{thm:error}, the relative error is $\approx 1.1 \times 10^{-14}$. Working in log-space, each of $\pi(10^6) = 78498$ terms incurs relative error $\varepsilon_{\mathrm{mach}} \approx 2^{-53}$, giving accumulated error $\sim 8.7 \times 10^{-12}$ in $L = \sum_p -\ln(1-p^{-3})$, which propagates to absolute error $\sim 10^{-11}$ in~$E$. This is well within the tolerance for the 15th decimal digit.
\end{proof}

\section{Algorithm}

\begin{enumerate}
\item \textbf{Sieve primes} up to $N = 10^6$ via the Sieve of Eratosthenes.
\item \textbf{Accumulate in log-space:} $L = \sum_{p \le N} \bigl[-\ln(1-p^{-3})\bigr]$.
\item \textbf{Recover:} $E(3, N) = e^L$.
\item \textbf{Extract:} $\lfloor E \times 10^{15} \rfloor$.
\end{enumerate}

\begin{verbatim}
function EulerProduct(N):
    is_prime[0..N] = all true
    is_prime[0] = is_prime[1] = false
    for i = 2 to floor(sqrt(N)):
        if is_prime[i]:
            for j = i*i to N step i:
                is_prime[j] = false
    L = 0.0
    for p = 2 to N:
        if is_prime[p]:
            L = L - log(1 - p^(-3))
    return floor(exp(L) * 10^15)
\end{verbatim}

\section{Complexity Analysis}

\begin{itemize}
\item \textbf{Time:} $O(N \log\log N)$ for the sieve, plus $O(\pi(N)) = O(N/\ln N)$ for log-space accumulation. Total: $O(N \log\log N)$.
\item \textbf{Space:} $O(N)$ for the sieve array.
\end{itemize}

\section{Answer}

\[
\boxed{495568995495726}
\]

\end{document}
